\documentclass{article}
\usepackage{ctex} 
\usepackage{amsmath, amssymb}
\usepackage{geometry}
\geometry{a4paper, margin=1in}

\title{深度强化学习知识点复习（价值学习部分）}
\author{本文档参照王树森深度强化学习进行补充整理}
\begin{document}
\maketitle

\section{训练神经网络总是用SGD而非GD原因}
\begin{itemize}
    \item 梯度下降（GD）在处理非凸问题时容易陷入鞍点，难以收敛到局部最优；而随机梯度下降（SGD）和小批量 SGD 由于梯度的随机性，有助于跳出鞍点，更可能趋近局部最优，且泛化能力更强。
    \item GD 每一步需要计算整个数据集上的梯度，计算量是 SGD 的 \(n\) 倍（\(n\) 为样本数量），通常非常慢，除非使用大规模并行计算；SGD 每次仅用一个或一批样本更新，计算效率高，适合在线学习。
\end{itemize}

\section{Robbins-Monro算法}
Robbins-Monro 算法是一种随机近似方法，用于在无法直接计算期望时在线更新参数。其更新形式为：
\[
\theta_t = \theta_{t-1} + \alpha_t \bigl( f(x_t) - \theta_{t-1} \bigr) = (1 - \alpha_t) \theta_{t-1} + \alpha_t  f(x_t) \tag{1}
\]
其中 \(\alpha_t\) 是步长序列。为保证均方收敛，\(\alpha_t\) 需满足：
\[
\sum_{t=1}^{\infty} \alpha_t = \infty, \quad \sum_{t=1}^{\infty} \alpha_t^2 < \infty \tag{2}
\]
该条件确保了算法能克服噪声并收敛到真实期望值，在 Q-learning 等强化学习算法中被用于更新动作价值函数。

\section{对比动作价值函数、最优动作价值函数和状态价值函数的依赖因素}
\begin{itemize}
    \item \textbf{动作价值函数} \(Q_\pi(s,a)\) 依赖于当前状态 \(s\)、当前动作 \(a\) 和策略 \(\pi\)，因为它衡量在策略 \(\pi\) 下从 \((s,a)\) 出发的期望回报。
    \item \textbf{最优动作价值函数} \(Q^*(s,a) = \max_\pi Q_\pi(s,a)\) 通过选择最优策略，只依赖于当前状态 \(s\) 和当前动作 \(a\)，与具体策略无关。
    \item \textbf{状态价值函数} \(V_\pi(s) = \mathbb{E}_{a\sim\pi(\cdot|s)}[Q_\pi(s,a)]\) 将动作视为随机变量求期望，因此只依赖于策略 \(\pi\) 和当前状态 \(s\)，不依赖于具体动作。
\end{itemize}

\section{价值学习中DQN算法训练框架}
DQN 使用时序差分（TD）算法学习最优动作价值函数，其理论基础是最优贝尔曼方程：
\[
Q^*(s,a) = \mathbb{E}_{s'\sim P(\cdot|s,a)} \left[ R + \gamma \max_{a'} Q^*(s',a') \right] \tag{3}
\]
训练时通过深度神经网络 \(Q(s,a;\theta)\) 近似 \(Q^*\)，并从经验回放池 \(\mathcal{D}\) 中采样转移 \((s,a,r,s')\)，最小化损失函数：
\[
L(\theta) = \mathbb{E}_{(s,a,r,s')\sim\mathcal{D}} \left[ \bigl( r + \gamma \max_{a'} Q(s',a';\theta^-) - Q(s,a;\theta) \bigr)^2 \right] \tag{4}
\]
其中 \(\theta^-\) 为目标网络参数（定期从 \(\theta\) 复制），用于稳定训练。DQN 是异策略（off-policy）方法，允许使用任意行为策略收集数据并反复利用。

\section{表格形式的Q学习训练框架}
表格 Q 学习直接基于最优贝尔曼方程进行在线更新。对于每个转移 \((s,a,r,s')\)，更新规则为：
\[
Q(s,a) \leftarrow Q(s,a) + \alpha \left[ r + \gamma \max_{a'} Q(s',a') - Q(s,a) \right] \tag{5}
\]
其中 \(\alpha\) 是学习率，\(\gamma\) 是折扣因子。该更新使用当前估计的 \(Q\) 值逼近最优贝尔曼方程中的期望，无需模型，属于异策略方法，因为更新中使用的 \(\max\) 操作隐含了目标策略为贪婪策略，而与行为策略无关。

\section{同策略on-policy与异策略off-policy方法比较}
在强化学习中，用于与环境交互收集经验的策略称为\textbf{行为策略}，而最终要学习优化的策略称为\textbf{目标策略}。若二者相同，则为同策略（如 SARSA）；若不同，则为异策略（如 Q-learning）。异策略的最大优势是可以通过经验回放池存储历史数据并反复使用，提高样本效率；同策略则要求每次更新必须使用当前策略新采集的数据，样本利用率较低，但通常收敛更稳定，方差更小。

\section{价值学习中表格形式的SARSA算法训练框架}
SARSA 基于贝尔曼方程对当前策略 \(\pi\) 下的动作价值函数进行估计：
\[
Q_\pi(s,a) = \mathbb{E}_{s'\sim P(\cdot|s,a),\,a'\sim\pi(\cdot|s')} \left[ R + \gamma Q_\pi(s',a') \right] \tag{6}
\]
对于每个五元组 \((s,a,r,s',a')\)（其中 \(a'\) 是根据当前策略 \(\pi\) 在 \(s'\) 下选择的动作），SARSA 采用以下在线更新：
\[
Q(s,a) \leftarrow Q(s,a) + \alpha \left[ r + \gamma Q(s',a') - Q(s,a) \right] \tag{7}
\]
SARSA 是同策略方法，因为更新中使用的 \(a'\) 必须由当前策略生成，若使用经验回放，则存储的 \(a'\) 可能来自旧策略，直接使用会导致更新方向偏离当前策略的贝尔曼方程，因此不能简单地应用经验回放。若需重用数据，必须结合重要性采样等技术，但基本 SARSA 不支持。

\section{价值学习中Q学习和SARSA的对比}
\begin{itemize}
    \item \textbf{更新公式}：Q学习使用 \(\max_{a'} Q(s',a')\)（公式 (5)），直接逼近最优价值函数；SARSA 使用当前策略下的动作 \(a'\) 对应的 \(Q(s',a')\)（公式 (7)），逼近当前策略的价值函数。
    \item \textbf{策略依赖性}：Q学习与行为策略无关（异策略），可使用任意策略收集的数据；SARSA 严格依赖于当前策略（同策略），更新必须使用当前策略选择的动作。
    \item \textbf{经验回放}：Q学习可配合经验回放，提高样本效率；SARSA 不能直接使用经验回放，因为旧经验中的 \(a'\) 不再反映当前策略。
    \item \textbf{收敛性}：在表格情形下，两者都能收敛到最优值函数（Q学习）或给定策略的值函数（SARSA），但 SARSA 的收敛通常更平稳，而 Q学习可能因最大化偏差导致过估计。
    \item \textbf{方差与偏差}：SARSA 使用实际采样的 \(a'\)，方差较小但可能陷入次优；Q学习使用最大化操作，偏差较小但方差较大，且容易引入正向偏差。
\end{itemize}

\section{价值学习中神经网络形式的SARSA算法}
当状态空间为无限集合时，可采用神经网络 \(Q(s,a;\theta)\) 来近似动作价值函数 \(Q_\pi(s,a)\)，其中 \(\theta\) 为网络参数。神经网络SARSA仍为同策略方法，需使用当前策略 \(\pi\) 采集五元组 \((s,a,r,s',a')\)，其中 \(a' \sim \pi(\cdot|s')\)。其训练目标是最小化TD误差的平方：
\[
L(\theta) = \mathbb{E}_{(s,a,r,s',a')\sim\mathcal{D}} \left[ \bigl( r + \gamma Q(s',a';\theta) - Q(s,a;\theta) \bigr)^2 \right] \tag{8}
\]
参数更新采用梯度下降：
\[
\theta \leftarrow \theta - \alpha \nabla_\theta L(\theta) \tag{9}
\]
注意，由于SARSA是同策略，经验回放池 \(\mathcal{D}\) 中存储的转移必须严格由当前策略生成，否则更新公式 (8) 不再对应当前策略的贝尔曼方程。在实际应用中，神经网络SARSA常与策略网络结合构成Actor-Critic方法：策略网络 \(\pi(a|s;\phi)\) 作为Actor控制智能体，价值网络 \(Q(s,a;\theta)\) 作为Critic评价当前策略并指导Actor更新。此时Critic仍采用SARSA方式更新，即基于实际采样的动作 \(a'\) 计算TD目标。

\section{多步TD算法的原理推导}
多步TD（n步TD）算法利用未来 \(n\) 步的累积奖励来更新价值函数，以平衡偏差与方差。其基于多步回报的贝尔曼方程：
\[
Q_\pi(s_t,a_t) = \mathbb{E}_{s_{t+1},a_{t+1},\dots,s_{t+n}} \left[ \sum_{k=0}^{n-1} \gamma^k r_{t+k} + \gamma^n Q_\pi(s_{t+n},a_{t+n}) \right] \tag{10}
\]
定义n步回报：
\[
G_t^{(n)} = \sum_{k=0}^{n-1} \gamma^k r_{t+k} + \gamma^n Q(s_{t+n},a_{t+n};\theta) \tag{11}
\]
则参数更新通过最小化TD误差实现：
\[
\theta \leftarrow \theta + \alpha \left( G_t^{(n)} - Q(s_t,a_t;\theta) \right) \nabla_\theta Q(s_t,a_t;\theta) \tag{12}
\]
训练流程：在每个时刻 \(t\)，智能体与环境交互并存储转移，当收集到 \(n\) 步数据后，利用公式 (11) 计算目标值，再按 (12) 更新网络。n步TD介于蒙特卡洛（\(n\to\infty\)）和一步TD（\(n=1\)）之间，可通过调整 \(n\) 控制偏差-方差权衡。

\section{训练价值网络采用蒙特卡洛方法的优劣性}
蒙特卡洛方法使用完整回报 \(G_t = \sum_{k=0}^{\infty} \gamma^k r_{t+k}\) 作为目标，其优势是\textbf{无偏估计}：由于回报直接来自真实轨迹，期望等于真实价值函数。但劣势是\textbf{方差大}，因为单条轨迹的随机性（如策略随机、环境转移随机）导致回报波动剧烈，且需等待回合结束才能更新，学习效率低。

\section{训练价值网络采用自举方法的优劣性}
自举方法（如一步TD）使用当前估计的后续价值构建目标，即 \(r + \gamma Q(s',a';\theta)\)。其优势是\textbf{方差小}，因为目标只依赖单步随机变量；劣势是\textbf{有偏估计}，因为目标中包含尚未收敛的当前价值估计，导致更新方向存在偏差。自举使得算法可在线更新，但可能引入累积误差。

\section{蒙特卡洛方法与自举的对比}
\begin{itemize}
    \item \textbf{偏差}：蒙特卡洛无偏；自举有偏（依赖于当前估计的准确性）。
    \item \textbf{方差}：蒙特卡洛高方差（需积分多步随机性）；自举低方差（仅涉及单步随机性）。
    \item \textbf{更新时机}：蒙特卡洛需回合结束；自举可每一步更新。
    \item \textbf{收敛性}：蒙特卡洛在表格情形下保证收敛，但收敛慢；自举可能加速收敛，但可能不收敛（尤其在函数近似时，如Q-learning可能发散）。
    \item \textbf{应用场景}：蒙特卡洛适用于回合制任务；自举适用于连续任务，常与函数近似结合。
\end{itemize}

\section{经验回放的优点}
\begin{enumerate}
    \item \textbf{打破序列相关性}：强化学习数据在时间上高度相关，直接使用连续样本更新违背了多数优化算法对独立同分布的假设。经验回放随机抽取历史转移，使得样本近似独立同分布，稳定训练。
    \item \textbf{提高样本效率}：每个转移可被多次使用，从有限交互中学习更多信息，减少与环境交互的总次数，尤其适用于数据获取成本高的场景。
\end{enumerate}

\section{样本数量与更新次数的区别}
\textbf{样本数量}指智能体与环境交互产生的转移总数（即采集的数据量），\textbf{更新次数}指利用这些样本对参数进行梯度下降的次数。在经验回放中，一个样本可被多次更新。样本数量更为关键，因为它决定了数据覆盖的状态-动作空间的广度；更新次数虽可增加，但过度重复使用有限样本可能导致过拟合和泛化能力下降。因此，收集更多样本来覆盖更广分布通常比单纯增加更新次数更有效。

\section{经验回放的局限性}
\begin{enumerate}
    \item \textbf{策略过时问题}：回放池中存储的转移由旧行为策略生成，与当前待学习的目标策略可能存在分布偏移。这可能导致更新方向偏离最优解，尤其在策略敏感算法（如同策略方法）中。
    \item \textbf{仅适用于异策略}：经验回放本质上要求目标策略与行为策略可以不同，因为存储的数据来自过去策略。同策略算法（如SARSA）若使用回放池，则存储的 \(a'\) 不再反映当前策略，直接更新会违背其理论基础。因此经验回放通常与异策略算法（如DQN、Q学习）结合使用。
\end{enumerate}

\section{优先经验回放原理及参数更新方式}
优先经验回放（PER）通过赋予经验样本不同的采样概率，使智能体更频繁地学习那些“意外性”高的样本，通常用 TD 误差衡量。其核心思想是：TD 误差绝对值越大，表明当前 Q 网络对该样本的预测越不准确，该样本应具有更高的采样优先级。

定义样本 \(i\) 的优先级 \(p_i > 0\)，常用两种形式：
\begin{itemize}
    \item 比例优先级：\(p_i = |\delta_i| + \epsilon\)，其中 \(\delta_i\) 为 TD 误差，\(\epsilon\) 是小常数防止概率为零。
    \item 基于排序的优先级：\(p_i = \frac{1}{\text{rank}(i)}\)，其中 \(\text{rank}(i)\) 是按 \(|\delta_i|\) 排序的序号。
\end{itemize}
采样概率由优先级归一化得到：
\[
P(i) = \frac{p_i^\alpha}{\sum_k p_k^\alpha}
\tag{13}
\]
其中 \(\alpha \in [0,1]\) 控制优先级的使用程度，\(\alpha=0\) 退化为均匀采样。

由于采样分布偏离原始经验分布，为无偏估计梯度，需乘以重要性采样权重：
\[
w_i = \left( \frac{1}{N} \cdot \frac{1}{P(i)} \right)^\beta
\]
其中 \(N\) 为回放缓冲区大小，\(\beta\) 从某个较小值（如 0.4）逐渐增至 1，以补偿偏差。实际更新时通常将权重归一化：\(w_i / \max_j w_j\) 以提高稳定性。

在 DQN 中，使用优先经验回放时，损失函数为：
\[
L(\theta) = \frac{1}{B} \sum_{i \in \text{batch}} w_i \left( r_i + \gamma \max_{a'} Q(s_i', a'; \theta^-) - Q(s_i, a_i; \theta) \right)^2
\tag{14}
\]
其中 \(B\) 为 batch 大小，\(\theta^-\) 为目标网络参数。每次更新后，需要重新计算被采样样本的 TD 误差，并更新其优先级 \(p_i\)。

\section{优先经验回放通过调整学习率修正}
优先经验回放引入的非均匀采样会改变损失函数的期望，导致有偏估计。重要性采样权重 \(w_i\) 的作用就是修正这一偏差：通过将每个样本的梯度乘以 \(w_i\)，使得更新方向近似于原始均匀采样下的期望梯度。

从参数更新角度看，这等价于对每个样本使用了自适应的学习率。设原始学习率为 \(\eta\)，则实际更新步长为 \(\eta w_i \nabla L_i\)。高优先级样本的 \(P(i)\) 大，对应的 \(w_i\) 小，从而降低了它的更新步长；低优先级样本的 \(w_i\) 大，步长相对提高。这样，虽然采样更频繁地偏向高优先级样本，但权重调整防止了这些样本过度主导更新，使整体更新更接近均匀采样的效果。

\section{大抽样概率、小学习率结合不会导致效果抵消原因分析}
高优先级样本被频繁采样（大抽样概率），但其权重 \(w_i\) 较小（小学习率），二者相乘会不会导致最终更新量不变，从而使优先采样的效果被抵消？

关键在于，优先采样改变的是样本被选中的频率，而权重调整的是每次更新的幅度。最终每个样本对参数更新的总贡献期望为：
\[
\mathbb{E}_{\text{PER}}[w_i \Delta_i] = \sum_i P(i) \cdot w_i \Delta_i = \sum_i P(i) \cdot \frac{1}{N P(i)} \Delta_i = \frac{1}{N} \sum_i \Delta_i = \mathbb{E}_{\text{uniform}}[\Delta_i]
\tag{15}
\]
其中 \(\Delta_i\) 为未加权的梯度。这说明，在期望意义上，PER 与均匀采样等价，即无偏性成立。然而，方差却降低了：高 TD 误差的样本贡献被分摊到更多更新步中，而低 TD 误差的样本贡献被集中，从而加速了学习。因此，大抽样概率和小学习率的结合并未抵消效果，反而在保持期望不变的前提下，通过减小方差提升了效率。

\section{Q 学习高估真实价值的原因}
Q 学习及其深度变种（如 DQN）常出现对动作价值的过高估计，主要原因有以下两点。
\begin{enumerate}
\item\textbf{自举导致偏差传播}
自举（bootstrapping）指用当前的估计值来更新自身。在 Q 学习中，更新目标为：
\[
Y_t = R_{t+1} + \gamma \max_{a} Q(S_{t+1}, a; \theta^-)
\tag{16}
\]
其中目标网络 \(\theta^-\) 本身也是估计值，包含误差。如果某次更新使 \(Q\) 被高估，那么后续使用该高估值作为目标时，误差会进一步传播放大。数学上，假设真实值 \(Q^*\)，估计值 \(\hat{Q}\)，初始误差 \(\epsilon_0\)，则经过一步更新：
\[
\hat{Q}_{\text{new}}(s,a) = \hat{Q}(s,a) + \alpha \left( r + \gamma \max_{a'} \hat{Q}(s',a') - \hat{Q}(s,a) \right)
\tag{17}
\]
若 \(\max_{a'} \hat{Q}(s',a') > \max_{a'} Q^*(s',a')\)，则目标偏高，导致 \(\hat{Q}_{\text{new}}\) 进一步高估。这种偏差通过贝尔曼方程迭代传播，最终可能导致严重的过高估计。

\item\textbf{最大化导致 TD 目标高估真实价值}
即使初始估计无偏，最大化操作也会引入正向偏差。设对于某个状态 \(s'\)，真实值为 \(Q^*(s',a)\)，而估计值 \(Q(s',a) = Q^*(s',a) + \varepsilon_a\)，其中 \(\varepsilon_a\) 是均值为 0 的独立噪声。则：
\[
\mathbb{E}\left[ \max_a Q(s',a) \right] \geq \max_a \mathbb{E}[Q(s',a)] = \max_a Q^*(s',a)
\tag{18}
\]
严格不等式成立通常因为最大值函数的期望大于期望的最大值。更具体地，假设所有动作的真实值相等（设为 \(V^*\)），噪声 \(\varepsilon_a\) 独立同分布，则：
\[
\mathbb{E}\left[ \max_a (V^* + \varepsilon_a) \right] = V^* + \mathbb{E}[\max_a \varepsilon_a] \geq V^*
\tag{19}
\]
且 \(\mathbb{E}[\max_a \varepsilon_a] > 0\)（只要分布非退化）。因此，即使噪声均值为零，取最大值后期望也会高估。
\end{enumerate}

\section{价值函数均匀高估与非均匀高估的影响}
若所有动作价值被均匀高估（即每个 \(Q(s,a)\) 都增加相同常数），则贪心策略不变，不会影响决策。但实际中高估是非均匀的：不同动作的噪声不同，且随着学习过程，某些动作可能被高估更多。这会导致：
\begin{itemize}
    \item 次优动作因高估而被选中，使策略变差。
    \item 高估在自举中传播，使状态价值整体膨胀，破坏学习的稳定性。
    \item 最终可能使算法无法收敛到最优策略。
\end{itemize}

\section{使用目标网络切断自筹误差传播}
自举（bootstrapping）是指用当前价值网络的估计值去更新自身，这会导致误差的传播和累积，使训练不稳定甚至发散。目标网络通过提供固定的TD目标来切断这一环路。

设在线网络参数为 $\theta$，目标网络参数为 $\theta^-$，两者结构相同但 $\theta^-$ 定期从 $\theta$ 复制（硬更新）或通过软更新逐步逼近。DQN的TD目标变为：
\[
Y_t = r_t + \gamma \max_{a'} Q(s_{t+1}, a'; \theta^-)
\tag{20}
\]
由于 $\theta^-$ 在 $C$ 步内保持不变，计算目标时使用的Q值不依赖于当前 $\theta$，从而避免了用刚刚更新的网络立即评估自身。这相当于将目标固定一段时间，使训练更稳定。

但需注意：目标网络只是延迟了自举，并未完全消除偏差，因为 $\theta^-$ 本身仍是旧参数的估计值，依然包含误差，只是更新频率降低。从理论上讲，目标网络使贝尔曼更新变成了一个半梯度（semi-gradient）方法，牺牲了一定无偏性换取了稳定性。

\section{双Q学习算法原理}
标准Q学习使用 $\max$ 操作同时进行动作选择和价值评估，容易导致过高估计。Double Q-learning 将选择和评估解耦，用两个独立的Q函数来减少正偏。

在Double DQN中，动作选择由在线网络 $\theta$ 完成，而价值评估由目标网络 $\theta^-$ 完成，目标为：
\[
Y_t = r_t + \gamma Q\big(s_{t+1}, \arg\max_{a'} Q(s_{t+1}, a'; \theta); \theta^- \big)
\tag{21}
\]
即先用 $\theta$ 选出最优动作 $a^* = \arg\max_{a'} Q(s_{t+1}, a'; \theta)$，再用 $\theta^-$ 评估该动作的价值。这样即使 $\theta$ 高估了某个动作，$\theta^-$ 的评估相对独立，可缓解高估。

参数更新仍通过最小化 TD 误差平方：
\[
L(\theta) = \mathbb{E}\left[ \big( Y_t - Q(s_t,a_t;\theta) \big)^2 \right]
\tag{22}
\]
理论上，Double DQN 能显著降低过高估计，使价值估计更准确，从而改善策略。

\section{最优优势函数}
优势函数 $A^\pi(s,a)$ 定义为动作价值与状态价值之差，表示采取动作 $a$ 相对于平均水平的优势：
\[
A^\pi(s,a) = Q^\pi(s,a) - V^\pi(s)
\tag{23}
\]
对于最优策略，记最优优势函数为 $A^*(s,a) = Q^*(s,a) - V^*(s)$。根据最优价值函数的定义，$V^*(s) = \max_a Q^*(s,a)$，代入可得：
\[
\max_a A^*(s,a) = \max_a Q^*(s,a) - V^*(s) = 0
\tag{24}
\]
即最优优势函数在最优动作上的值为0。这一性质为后续对决网络的设计提供了理论基础：可以将Q值分解为状态价值加上优势函数，并通过对优势函数施加约束来保证分解的唯一性。

\section{对决网络结构与原理}
对决网络（Dueling Network）是对 DQN 架构的改进，它将 Q 值显式分解为状态价值 $V(s)$ 和优势函数 $A(s,a)$ 两部分：
\[
Q(s,a; \theta, \alpha, \beta) = V(s; \theta, \beta) + A(s,a; \theta, \alpha)
\tag{25}
\]
其中 $\theta$ 是卷积层等共享参数，$\beta$ 和 $\alpha$ 分别为价值流和优势流的参数。这种结构使网络能分别学习状态价值和动作优势，在某些情况下（如某些动作对价值影响不大时）收敛更快、估计更准。

其原理是：在训练过程中，价值流会学习状态的基准价值，而优势流则学习每个动作相对于基准的偏移。通过共享底层特征，两部分相互促进，提高了表示效率。实验表明，对决网络在同样参数规模下优于传统 DQN。

\section{解决不唯一性，max D(s,a,w)不能省略}
上述分解存在\textbf{不唯一性}：若将 $V$ 加上任意常数 $c$，同时将 $A$ 减去 $c$，则 $Q$ 保持不变。这导致网络无法确定 $V$ 和 $A$ 的取值，可能产生无意义的波动，影响收敛。

根据 $\max_a A^*(s,a)=0$ 的性质，可以强制优势函数在最优动作处为零，从而消除不唯一性。一种常用做法是：
\[
Q(s,a; \theta, \alpha, \beta) = V(s; \theta, \beta) + \left( A(s,a; \theta, \alpha) - \max_{a'} A(s,a'; \theta, \alpha) \right)
\tag{26}
\]
这样 $\max_a Q(s,a) = V(s)$，保证了 $V$ 确实是状态价值，$A$ 唯一地表示优势。此处的 $\max$ 项不能省略，否则无法施加约束，分解仍然不唯一。

\section{实际实现对决网络时为什么用mean替代max}
尽管理论公式使用 $\max$ 操作，但实践中通常用均值替代最大值：
\[
Q(s,a) = V(s) + \left( A(s,a) - \frac{1}{|\mathcal{A}|} \sum_{a'} A(s,a') \right)
\tag{27}
\]
原因如下：
\begin{itemize}
    \item \textbf{梯度稳定性}：$\max$ 操作会使梯度只更新被选中的动作对应的优势，其他动作的优势梯度为0，导致更新不稳定；而均值使所有动作的优势都参与梯度更新，更加平滑。
    \item \textbf{收敛速度}：均值方式下，优势函数倾向于更快地收敛到相对值，且实验表明整体收敛更快。
    \item \textbf{等价性}：虽然此时 $V(s)$ 不再严格等于 $\max_a Q(s,a)$，但只是整体偏移了一个常数，不影响动作选择的相对顺序，因此对最终策略无影响。
\end{itemize}

\section{噪声网络原理}
噪声网络（NoisyNet）是一种替代 $\epsilon$-greedy 的探索策略，通过在网络参数上添加参数化噪声来驱动探索，使探索成为网络的一部分并随训练自适应调整。

对于线性层 $y = wx + b$，噪声网络将其参数化为：
\[
w = \mu^w + \sigma^w \odot \varepsilon^w, \quad b = \mu^b + \sigma^b \odot \varepsilon^b
\tag{28}
\]
其中 $\mu$ 和 $\sigma$ 是可学习参数，$\varepsilon$ 是零均值噪声（如标准正态分布），$\odot$ 为逐元素乘。前向传播时，采样 $\varepsilon$ 得到实际参数，从而输出的Q值带有随机性，驱动探索。噪声的方差由 $\sigma$ 控制，网络可通过调整 $\sigma$ 自适应地控制探索程度：初期 $\sigma$ 较大，探索强；后期 $\sigma$ 可能减小，探索减弱。

梯度反向传播时，只更新 $\mu$ 和 $\sigma$，$\varepsilon$ 不参与更新。噪声的引入使探索与网络融为一体，避免了手工调参 $\epsilon$，且在理论上表现更好。

\section{噪声 DQN 实现框架}
噪声 DQN 将原始 DQN 中的全连接层替换为噪声层，实现步骤如下：
\begin{enumerate}
    \item 对每个要替换的线性层，定义可学习参数 $\mu \in \mathbb{R}^{d_{\text{out}}\times d_{\text{in}}}$ 和 $\sigma \in \mathbb{R}^{d_{\text{out}}\times d_{\text{in}}}$（偏置同理）。
    \item 每次前向传播时，采样噪声 $\varepsilon^w \sim \mathcal{N}(0,1)$ 和 $\varepsilon^b \sim \mathcal{N}(0,1)$，计算实际权重和偏置：
    \[
    w = \mu^w + \sigma^w \odot \varepsilon^w,\quad b = \mu^b + \sigma^b \odot \varepsilon^b
    \]
    \item 使用噪声参数进行前向计算得到Q值。
    \item 训练时，损失函数与 DQN 相同（如公式(4)），梯度反向传播更新 $\mu$ 和 $\sigma$。
    \item 探索由参数的随机性自动产生，无需 $\epsilon$-greedy。
\end{enumerate}
实践中，为减少计算量，常采用\textbf{因子化噪声}：每个单元独立采样一个噪声，再通过线性变换构造参数噪声，避免为每个参数独立采样。





\end{document}